% ---------------------------------------------------------------
% Talk: p-adic Variational Quantum Circuits Have No Barren Plateaus
% AIQxQIA 2026, Perugia -- Thursday 8 October 2026, 11:55-12:15
% Session "AI for Quantum - 2", full paper: 17 min + 3 min questions.
% Author: Piero Giacomelli
%
% Compile: pdflatex Paper_14_Giacomelli_piero && pdflatex Paper_14_Giacomelli_piero
% Upload as: Paper_14_Piero_Giacomelli.pdf
% ---------------------------------------------------------------
\documentclass[aspectratio=169,compress]{beamer}
\usetheme{Copenhagen}
\usecolortheme[RGB={31,59,112}]{structure}

\usepackage{amsmath,amssymb}
\usepackage{booktabs}
\usepackage{tikz}

\setbeamertemplate{navigation symbols}{}

% ---- notation, matching the paper ----
\newcommand{\Qp}{\mathbb{Q}_p}
\newcommand{\Zp}{\mathbb{Z}_p}
\newcommand{\OK}{\mathcal{O}}
\newcommand{\mm}{\mathfrak{m}}
\newcommand{\FF}{\mathbb{F}}
\newcommand{\Uni}{\mathrm{U}}
\newcommand{\Id}{\mathbf{1}}
\newcommand{\loss}{\mathcal{L}}
\newcommand{\Sph}{\mathbb{S}}
\newcommand{\Nm}{\mathrm{Nm}}
\newcommand{\padicnorm}[1]{\left|#1\right|_p}
\newcommand{\inner}[2]{\langle #1, #2\rangle}

\title[No Barren Plateaus over $\Qp$]{$p$-adic Variational Quantum Circuits\\
  Have No Barren Plateaus}
\author[P. Giacomelli]{Piero Giacomelli}
\institute[Tenax SPA]{Tenax SPA, Verona, Italy}
\date[AIQxQIA 2026]{AIQxQIA 2026 -- Perugia, 8 October 2026}

\begin{document}

\begin{frame}
  \titlepage
\end{frame}

\begin{frame}{Agenda}
  \tableofcontents
\end{frame}

% =================================================================
\section[Obstruction]{The obstruction}

\begin{frame}{Barren plateaus in variational quantum circuits}
  A variational circuit prepares $U(\theta)\psi_0$ and tunes $\theta$ to
  minimise
  \[
    \loss(\theta)=\inner{U(\theta)\psi_0}{M\,U(\theta)\psi_0}.
  \]
  \begin{block}{McClean et al.\ (2018)}
    If the circuit ensemble forms a unitary $2$-design on $\Uni(2^{n})$, then
    \[
      \mathbb{E}[\partial_\theta\loss]=0,
      \qquad
      \mathrm{Var}[\partial_\theta\loss]=O(2^{-n}).
    \]
  \end{block}
  \begin{itemize}
    \item The landscape is flat almost everywhere.
    \item Resolving a descent direction costs a number of shots exponential in
      the qubit count.
  \end{itemize}
\end{frame}

\begin{frame}{Why it matters for quantum machine learning}
  \begin{itemize}
    \item Barren plateaus are the main argument against \alert{scalable}
      variational quantum machine learning.
    \item Later work traced them to the cost function (Cerezo et al.\ 2021) and
      to ansatz expressibility (Holmes et al.\ 2022).
    \item The known mitigations act on the \alert{circuit}:
      \begin{itemize}
        \item initialisation and warm starts,
        \item local cost functions,
        \item shallow or structured ans\"atze.
      \end{itemize}
  \end{itemize}
  \vspace{0.6em}
  \begin{center}
    Every one of them keeps the scalars: $\mathbb{C}$.
  \end{center}
\end{frame}

\begin{frame}{What does the theorem actually use?}
  Read carefully, the barren-plateau theorem rests on three ingredients:
  \vspace{0.4em}
  \begin{enumerate}
    \item a \alert{compact} ensemble, so a Haar \emph{probability} measure
      exists;
    \item second-moment data supplied by unitary \alert{$t$-designs};
    \item ``small'' measured in the \alert{Archimedean} absolute value.
  \end{enumerate}
  \vspace{0.8em}
  None of the three is a fact about quantum mechanics.
  All three are facts about $\mathbb{C}$.
  \vspace{0.8em}
  \begin{center}
    \Large\emph{What happens if we change the field?}
  \end{center}
\end{frame}

% =================================================================
\section[Over $\Qp$]{Quantum computing over $\Qp$}

\begin{frame}{Changing the scalars}
  Following Aniello, Mancini and Parisi, rebuild quantum kinematics over the
  unramified quadratic extension $K=\Qp(\sqrt\mu)$ of the $p$-adic field:
  \begin{itemize}
    \item states are vectors in $K^{N}$, with $N=2^{n}$;
    \item gates are unitaries for the Hermitian form
      $\inner vw=\sum_i\bar v_iw_i$;
    \item the geometry is \alert{ultrametric}:
      $\padicnorm{x+y}\le\max\{\padicnorm{x},\padicnorm{y}\}$.
  \end{itemize}
  \begin{block}{Why anyone builds this}
    Ultrametric geometry is the intrinsic geometry of \alert{hierarchical
    data}: taxonomies, phylogenies, file systems. Classical $p$-adic learning
    has used it for decades.
  \end{block}
  The quantum side already has a qubit, a tensor product, a universal gate set
  and a first quantum neural network.
\end{frame}

\begin{frame}{Fact 1: small means divisible}
  \begin{columns}[T]
    \begin{column}{0.40\textwidth}
      $\padicnorm{x}=p^{-v_p(x)}$, where $v_p(x)$ counts the factors of $p$
      in $x$.
      \vspace{0.6em}

      At $p=3$:
      \[
        |18|_3=\tfrac19\ \ (18=2\cdot3^{2}),\qquad
        |2|_3=1 .
      \]
      Two numbers are close when they agree \alert{modulo a high power of
      $p$}: they sit deep in the same branch of the tree.
      \vspace{0.6em}

      Distances are read off a hierarchy: this is why the geometry fits
      hierarchical data.
    \end{column}
    \begin{column}{0.58\textwidth}
      \centering
      \begin{tikzpicture}[
          every node/.style={font=\scriptsize},
          dot/.style={circle,fill=structure.fg,inner sep=1.4pt},
          hot/.style={circle,fill=red!75!black,inner sep=1.8pt},
          level distance=11mm]
        \node[dot,label=above:{$\mathbb{Z}_3$}] (r) at (0,0) {};
        % level 1: residues mod 3
        \node[hot,label=left:{$0$}] (a0) at (-2.2,-1.1) {};
        \node[dot,label=left:{$1$}] (a1) at (0,-1.1) {};
        \node[dot,label=right:{$2$}] (a2) at (2.2,-1.1) {};
        \draw (r)--(a0) (r)--(a1) (r)--(a2);
        % level 2 under 0: residues mod 9
        \node[hot,label=left:{$0$}] (b0) at (-3.0,-2.2) {};
        \node[dot,label=below:{$3$}] (b3) at (-2.2,-2.2) {};
        \node[dot,label=below:{$6$}] (b6) at (-1.4,-2.2) {};
        \draw (a0)--(b0) (a0)--(b3) (a0)--(b6);
        % level 2 under 2: residues mod 9
        \node[dot,label=below:{$2$}] (c2) at (1.4,-2.2) {};
        \node[dot,label=below:{$5$}] (c5) at (2.2,-2.2) {};
        \node[dot,label=below:{$8$}] (c8) at (3.0,-2.2) {};
        \draw (a2)--(c2) (a2)--(c5) (a2)--(c8);
        % level 3 under 0 mod 9: residues mod 27
        \node[dot,label=below:{$0$}] (d0) at (-3.7,-3.3) {};
        \node[dot,label=below:{$9$}] (d9) at (-3.0,-3.3) {};
        \node[hot,label=below:{$18$}] (d18) at (-2.3,-3.3) {};
        \draw (b0)--(d0) (b0)--(d9) (b0)--(d18);
        \draw[red!75!black,very thick] (r)--(a0)--(b0)--(d18);
        % level labels
        \node[anchor=west,gray] at (3.3,-1.1) {mod $3$};
        \node[anchor=west,gray] at (3.3,-2.2) {mod $9$};
        \node[anchor=west,gray] at (3.3,-3.3) {mod $27$};
        \node[anchor=west,text width=2.6cm,red!75!black] at (-1.6,-3.4)
          {$18$: two levels in the $0$-branch, so $|18|_3=3^{-2}$};
      \end{tikzpicture}
    \end{column}
  \end{columns}
\end{frame}

\begin{frame}{Facts 2 and 3: a discrete scale, and sums that do not grow}
  \begin{block}{2. The scale is discrete}
    $\padicnorm{\cdot}$ takes values in $p^{\mathbb{Z}}$. A gradient is a unit,
    or divisible by $p$, or by $p^{2}$, and nothing in between. There is no
    continuum of ``slightly small'' quantities.
  \end{block}
  \begin{block}{3. Sums do not grow; cancellation is a congruence}
    $\padicnorm{x+y}=\max\{\padicnorm{x},\padicnorm{y}\}$ whenever the two
    sizes differ. A sum is as large as its largest term \emph{unless} the
    leading digits cancel, and that is a condition modulo $p$: finite and
    checkable.
  \end{block}
  \vspace{0.3em}
  Fact 3 is why the whole second-moment analysis will collapse onto
  \alert{counting in finite groups}.
\end{frame}

\begin{frame}{One thing that does not carry over: sampling}
  \begin{itemize}
    \item $K$ is not ordered, so $\inner\psi\psi\ge0$ means nothing: states
      form an affine, not a convex, set.
    \item Generalised probabilities lie in $\Qp$, not in $[0,1]$. They have no
      frequentist reading and \alert{cannot be estimated by repeated
      measurement}.
    \item Readout is algebraic: the algorithm returns $p$-adic numbers, read
      digit by digit.
  \end{itemize}
  \begin{alertblock}{So what is a plateau here?}
    Over $\mathbb{C}$ a small gradient is bad because it takes many shots to
    resolve. Over $K$ there are no shots: a gradient of high valuation is bad
    because the \emph{leading digits of the parameter update vanish}.
  \end{alertblock}
\end{frame}

% =================================================================
\section[Question]{Asking the question correctly}

\begin{frame}{Surprise 1: the usual ensemble does not exist}
  \begin{alertblock}{Proposition}
    For every $N\ge2$ the unitary group $\Uni(K^{N})$ is unbounded, hence
    non-compact, and carries \emph{no} Haar probability measure.
  \end{alertblock}
  \begin{itemize}
    \item Explicit witness at $p=3$, $K=\mathbb{Q}_3(i)$: the vector
      $u=(\tfrac13,c)$ with $\Nm(c)=\tfrac89$ has $\inner uu=1$ but
      $|\tfrac13|_3=3$. Witt's theorem completes it to a unitary of norm $3$;
      replacing $\tfrac13$ by $3^{-m}$ gives norm $3^{m}$.
    \item The reason is structural: over a local field $\sum\bar z_iz_i$ is
      \alert{isotropic}. Over $\mathbb{C}$ compactness came from
      positivity, and positivity is what the $p$-adic model gives up.
  \end{itemize}
  \vspace{0.3em}
  ``Haar-random circuit'' over the full group has no meaning.
\end{frame}

\begin{frame}{The fix: integral unitaries}
  \begin{block}{Definition}
    $G_N=\Uni_N(\OK)=\{U\in\Uni(K^{N}) : \text{all entries of $U$ lie in the
    ring of integers } \OK\}$.
  \end{block}
  \begin{itemize}
    \item $G_N$ is compact and carries a unique Haar probability measure $\nu$.
    \item It costs nothing for quantum computing: it contains every circuit
      built from the known $p$-adic gate sets:
      \begin{itemize}
        \item Pauli gates and CNOT (entries in $\{0,\pm1\}$);
        \item the $\mathrm{SO}(3)_p$ gates of Svampa et al.\ (roots of unity);
        \item the Hadamard gate, whenever $\sqrt2\in K$.
      \end{itemize}
  \end{itemize}
  \vspace{0.3em}
  $G_N$ is the smallest natural compact group containing all realisable
  circuits, so it is our ensemble.
\end{frame}

\begin{frame}{Surprise 2: the benchmark is $p^{-1}$, not $1$}
  The series $\exp_p(X)=\sum_m X^{m}/m!$ converges only for
  $\|X\|<p^{-1/(p-1)}$. Because the value group is discrete, this is
  \alert{equivalent} to $\|X\|\le p^{-1}$, that is $X\equiv0\bmod p$.
  \vspace{0.5em}

  Hence every convergent one-parameter ansatz satisfies
  \[
    \exp_p(\theta X)\equiv\Id\pmod p
    \qquad\Longrightarrow\qquad
    \padicnorm{\partial_\theta\loss}\le p^{-1}\ \text{ always.}
  \]
  \begin{itemize}
    \item We therefore write the generator as $p\,\theta A$ with $A$ integral:
      the factor $p$ is not a normalisation we chose, it \emph{is} the
      convergence condition.
    \item ``As large as possible'' means $\padicnorm{\partial_\theta\loss}
      =p^{-1}$, i.e.\ the normalised gradient
      $\nabla\loss=p^{-1}\partial_\theta\loss$ is a $p$-adic unit.
  \end{itemize}
\end{frame}

\begin{frame}{The question, stated precisely}
  $\mathrm{Var}=O(2^{-n})$ has no $p$-adic meaning. The faithful translation
  asks how often the gradient is \emph{as large as it can be}.
  \begin{block}{Question}
    Let $\nu$ be Haar measure on $G_N$, $N=2^{n}$. Fix an integral observable
    $M=M^*$, an integral generator $A=-A^*$, an admissible $\psi_0$ and a block
    $U_B\in G_N$, and set
    \[
      U(\theta)=U_B\,\exp_p(p\,\theta A)\,U_A,
      \qquad \theta\in\Zp .
    \]
    If $U_A\sim\nu$, does
    $\Pr\bigl[\padicnorm{\partial_\theta\loss}=p^{-1}\bigr]$ stay bounded away
    from $0$ as $n\to\infty$, and for which pairs $(M,A)$?
  \end{block}
  \textbf{Heuristic.} A Haar-random element of $\Zp$ is a unit with
  probability $1-\frac1p$. So one expects the answer $\approx1-\frac1p$,
  whatever $n$.
\end{frame}

% =================================================================
\section[Result]{Main result}

\begin{frame}{No barren plateaus}
  Let $\widehat M=U_B^{*}MU_B$, and let $\Gamma$ be the commutator
  $[\widehat M,A]$ reduced modulo $p$.
  \begin{block}{Definition ($r$-non-degenerate)}
    $\Gamma$ has an orthonormal eigenbasis over $\FF_{p^2}$ with eigenvalues in
    $\FF_p$ (\emph{splitting}), and no eigenvalue has multiplicity above $N-r$.
  \end{block}
  \begin{alertblock}{Theorem}
    If $\Gamma$ is $r$-non-degenerate and $U_A$ is drawn from $\nu$, then
    \[
      \Pr\Bigl[\padicnorm{\partial_\theta\loss}=p^{-1}\Bigr]
      \;\ge\;1-\frac1p-10\,p^{\,1-r}.
    \]
  \end{alertblock}
  The gradient attains the largest absolute value any convergent $p$-adic
  ansatz permits, with probability bounded below \alert{independently of the
  number of qubits}. The same bound holds if $U_B$ is Haar-random too.
\end{frame}

\begin{frame}{Reading the theorem}
  \begin{itemize}
    \item \textbf{No dependence on $n$.} Over $\mathbb{C}$, Chebyshev gives
      $\Pr[|\partial_\theta\loss|\ge\varepsilon]=O(2^{-n}\varepsilon^{-2})$.
      Over $K$ the gradient is \emph{maximal} with probability close to
      $1-\frac1p$, for every $n$.
    \item \textbf{$r$ grows like $2^{n}$.} $\Gamma$ has $N=2^{n}$ eigenvalues
      in a $p$-element field; spread over the available values,
      $r=\Theta(2^{n})$.
    \item \textbf{Doubly exponential correction.} The error term $10p^{1-r}$ is
      of order $p^{-2^{n}}$, smallest exactly where the complex plateau is
      worst: many qubits, expressive circuits.
    \item \textbf{When it bites.} The bound is non-trivial once
      $p^{r-1}>10p/(p-1)$, which in the regime $r\sim2^{n}$ holds from three
      qubits on.
  \end{itemize}
\end{frame}

\begin{frame}{A witness you would actually write down}
  Take $p\equiv3\pmod4$, so $K=\Qp(i)$. On $n$ qubits:
  \[
    M=X_1,
    \qquad
    A=-\tfrac{i}{2}\bigl(\Id-Z_1\bigr),
    \qquad U_B=\Id .
  \]
  In words: \emph{measure Pauli $X$ on the first qubit, rotate the first
  qubit about $Z$.}
  \vspace{0.4em}
  \begin{itemize}
    \item $C=[M,A]=Y_1$, so $\Gamma^{2}=\Id$: eigenvalues $\pm1\in\FF_p$, each
      of multiplicity $2^{n-1}$, with orthonormal eigenbases.
    \item Hence $\Gamma$ is $r$-non-degenerate with $r=2^{\,n-1}$.
  \end{itemize}
  \begin{exampleblock}{At $n=10$, $p=3$}
    \[
      \Pr\bigl[\text{gradient maximal}\bigr]\;\ge\;1-\tfrac13-3^{-500}.
    \]
  \end{exampleblock}
\end{frame}

% =================================================================
\section[Proof]{Proof sketch}

\begin{frame}{Step 1: the gradient is a Hermitian form}
  Write $\psi=U_A\psi_0$. Expanding $\exp_p(p\theta A)$ to first order and
  using $A^{*}=-A$:
  \[
    \nabla\loss\;=\;p^{-1}\partial_\theta\loss
    \;=\;\inner{\psi}{C\,\psi},
    \qquad
    C=[\widehat M,A]=\widehat MA-A\widehat M .
  \]
  \begin{itemize}
    \item $C$ is selfadjoint, so $\nabla\loss\in\Qp$; all data integral gives
      $\nabla\loss\in\Zp$.
    \item The gradient is maximal $\iff$ $\nabla\loss$ is a unit $\iff$ its
      \alert{leading digit}
      \[
        \inner{\tilde\psi}{\Gamma\,\tilde\psi}\in\FF_p
      \]
      is non-zero, where $\tilde{\ }$ denotes reduction modulo $p$.
  \end{itemize}
  \vspace{0.3em}
  The question is now about one Hermitian form over a \emph{finite} field,
  evaluated at a random point.
\end{frame}

\begin{frame}{Step 2: Haar measure becomes counting}
  Compactness of $G_N$ comes with a filtration by reduction:
  \[
    G_N\;\xrightarrow{\ \pi_k\ }\;\Uni_N(\OK/\mm^{k})
    \;\longrightarrow\;\cdots\;\longrightarrow\;
    \Uni_N(\FF_{p^2}).
  \]
  \begin{block}{Pushforward}
    Each $\pi_k$ is surjective (lift, then correct by Newton iteration), and
    sends $\nu$ to the \alert{uniform} measure on the finite group
    $\Uni_N(\OK/\mm^{k})$.
  \end{block}
  \begin{block}{Transitivity}
    $G_N$ acts transitively on the integral sphere
    $\{u:\inner uu=1\}$, so $\tilde\psi$ is \alert{uniform} on the finite
    sphere $\widetilde\Sph_N=\{\tilde u\in\FF_{p^2}^{N}:\inner{\tilde u}{\tilde
    u}=1\}$.
  \end{block}
  This replaces the $t$-design combinatorics of the complex theory.
\end{frame}

\begin{frame}{Step 3: a character sum in closed form}
  In an orthonormal eigenbasis of $\Gamma$, with $t_i=\Nm(u_i)\in\FF_p$, the
  two constraints are $\sum_i t_i=1$ and $\sum_i\lambda_it_i=c$. Count the
  solutions $\mathcal{N}_c$ with an additive character $\chi$ of $\FF_p$.
  Everything reduces to the sum
  $T(s)=\sum_{t\in\FF_p}\#\Nm^{-1}(t)\,\chi(st)$, which is
  \alert{$p^{2}$ for $s=0$ and $-p$ for $s\neq0$}.
  \begin{alertblock}{Counting lemma}
    If $\Gamma$ is $r$-non-degenerate, then for every $c\in\FF_p$
    \[
      \Bigl|\frac{\mathcal{N}_c}{|\widetilde\Sph_N|}-\frac1p\Bigr|\le10\,p^{\,1-r},
      \qquad |\widetilde\Sph_N|=p^{2N-1}-(-1)^{N}p^{N-1}.
    \]
  \end{alertblock}
  Take $c=0$ (the leading digit vanishes) and the complement: that is the
  theorem.
\end{frame}

% =================================================================
\section{Discussion}

\begin{frame}{Side by side}
  \begin{center}
  \small
  \begin{tabular}{@{}lll@{}}
    \toprule
     & \textbf{Complex} & \textbf{$p$-adic} \\
    \midrule
    Ensemble        & $\Uni(2^{n})$, compact       & full group non-compact; use $G_N$ \\
    Moments         & unitary $t$-designs          & character sums over $\Uni_N(\FF_{p^2})$ \\
    Gradient        & $\mathrm{Var}=O(2^{-n})$     & $\Pr[\text{maximal}]\ge1-p^{-1}-10p^{1-r}$ \\
    Scaling in $n$  & exponential decay            & none \\
    Bad set         & full measure as $n\to\infty$ & locus $[\widehat M,A]\equiv0 \bmod p$ \\
    Expressibility  & more expressive, flatter     & more random, \alert{larger} gradients \\
    \bottomrule
  \end{tabular}
  \end{center}
  \vspace{0.6em}
  The last row is the one to remember: over $K$, randomness is what
  \emph{produces} large gradients.
\end{frame}

\begin{frame}{What replaces the plateau}
  The obstruction does not vanish: it changes shape.
  \begin{itemize}
    \item \textbf{Complex:} a \emph{measure} phenomenon. The bad set has full
      measure as $n\to\infty$; no choice of $M$ and $A$ escapes it.
    \item \textbf{$p$-adic:} a \emph{locus}. The bad set is exactly
      $[\widehat M,A]\equiv0\bmod p$: reduce two integer matrices and multiply
      them.
  \end{itemize}
  \begin{itemize}
    \item As $U_B$ varies, the locus is a clopen union of congruence classes of
      measure $O(p^{-1})$: small but \alert{not null}, a distinction the
      complex theory never has to make.
  \end{itemize}
  \begin{block}{Design criterion}
    Choose $\widehat M$ and $A$ that visibly fail to commute modulo $p$. The
    most ordinary pair (Pauli $X$ against a phase rotation) already does.
  \end{block}
\end{frame}

\begin{frame}{Does a maximal gradient imply trainability?}
  Not by itself.
  \begin{itemize}
    \item The theorem says the \emph{leading digit} of the gradient is
      generically non-zero, so a step $\theta\leftarrow\theta-\eta\nabla\loss$
      with $\padicnorm\eta=1$ moves the leading digit of $\theta$.
    \item Convergence is a separate condition. On a quadratic loss with Hessian
      $H$, $p$-adic descent contracts only if $\Id-\eta H\equiv0\bmod p$:
      with a scalar step, $H$ must be \alert{scalar modulo $p$}.
    \item With a matrix preconditioner the condition relaxes to
      \alert{invertibility of $H$ modulo $p$}, which is generic.
    \item When it holds, convergence is ultra-fast: error contracts by $p^{-1}$
      per step, so $O(\log_p(1/\varepsilon))$ iterations.
  \end{itemize}
  \vspace{0.3em}
  We remove gradient concentration from the list of obstructions; trainability
  becomes a question of \emph{conditioning}.
\end{frame}

\begin{frame}{Why study a model with no hardware?}
  No $p$-adic quantum computer exists. The result is still useful.
  \begin{enumerate}
    \item \alert{A diagnostic about the complex theory.} We keep the circuit,
      observable, ansatz and qubit count, change only the field, and the
      plateau disappears. The load is carried by compactness and by the
      Archimedean absolute value, which locates the hypothesis a mitigation has
      to attack.
    \item \alert{A reusable method.} The pushforward reduces \emph{any} Haar
      moment over $G_N$ to counting in finite unitary groups: overlaps,
      expectation values, $p$-adic kernels.
    \item \alert{A checkable criterion.} $[\widehat M,A]\not\equiv0\bmod p$ is
      a finite test: no sampling, no asymptotics.
  \end{enumerate}
  Applications would come through classical $p$-adic learning on hierarchical
  data, where the geometry is ultrametric to begin with.
\end{frame}

\begin{frame}{Limitations}
  \begin{itemize}
    \item \textbf{The splitting hypothesis is real.} Over a finite field a
      Hermitian matrix may have eigenvalues outside $\FF_p$; we give a
      $2\times2$ counterexample over $\FF_9$. Outside the splitting class the
      question is open.
    \item \textbf{Level one only.} We control $\Pr[v_p(\nabla\loss)\ge1]$, not
      the full valuation distribution (expected:
      $\Pr[v_p\ge k]=p^{-k}(1+o(1))$, not proved).
    \item \textbf{Standing assumptions.} $p$ odd, $K/\Qp$ unramified, one
      parameter in a two-block ansatz.
    \item \textbf{No experiments.} No sampling-based readout exists. The
      useful companion is exact arithmetic in $\OK/\mm^{k}$ for small $n$ and
      $p$, checking the counting lemma by enumeration.
  \end{itemize}
\end{frame}

\begin{frame}{Closing}
  \begin{center}
    {\Large Barren plateaus are a theorem about the\\[0.2em]
    Archimedean absolute value.}\\[1em]
    {\large Change the field and they go.}\\[2em]
    Piero Giacomelli \,$|$\, Tenax SPA \,$|$\, pgiacome@gmail.com\\[0.4em]
    \small Paper 14, AIQxQIA 2026 (CEUR-WS proceedings)
  \end{center}
\end{frame}

% =================================================================
\appendix

\begin{frame}{Backup: why the splitting hypothesis is not free}
  Over $\FF_9=\FF_3(\iota)$, $\iota^{2}=-1$, take
  \[
    \Gamma=\begin{pmatrix}0 & 1+\iota\\ 1-\iota & 0\end{pmatrix},
    \qquad
    \chi_\Gamma(x)=x^{2}-\Nm(1+\iota)=x^{2}+1 .
  \]
  \begin{itemize}
    \item $\Gamma$ is Hermitian, but its eigenvalues $\pm\iota$ lie in
      $\FF_9\setminus\FF_3$.
    \item If $\Gamma v=\lambda v$ with $\lambda\ne\bar\lambda$, then
      $\lambda\inner vv=\bar\lambda\inner vv$ forces $\inner vv=0$: the
      eigenvector $(1+\iota,\iota)$ has norm $2+1=0$.
    \item No orthonormal eigenbasis exists.
  \end{itemize}
  Closing the gap needs the classification of \emph{pencils} of Hermitian forms
  over a finite field: the main open follow-up.
\end{frame}

\begin{frame}{Backup: how large can $r$ be?}
  $\Gamma$ has $N=2^{n}$ eigenvalues in the $p$-element field $\FF_p$, so
  pigeonhole gives $m_{\max}\ge\lceil N/p\rceil$ and
  \[
    r\;\le\;N-\lceil N/p\rceil\;\le\;N\Bigl(1-\tfrac1p\Bigr).
  \]
  \begin{itemize}
    \item A simple spectrum is impossible once $N>p$.
    \item Spread evenly, $m_{\max}\approx N/p$ and $r\approx N(1-1/p)$.
    \item The witness spreads them over two values: $r=N/2$.
  \end{itemize}
  Either way $r=\Theta(2^{n})$, so $10p^{1-r}$ is of order $p^{-2^{n}}$.
\end{frame}

\begin{frame}{Backup: the bad set is a commutator}
  \begin{block}{Corollary}
    If $\Gamma$ splits, it fails to be $1$-non-degenerate if and only if
    $\widehat M$ and $A$ commute modulo $p$.
  \end{block}
  \begin{itemize}
    \item Failure means $m_{\max}=N$, i.e.\ $\Gamma=\lambda\Id$.
    \item A commutator has trace $0$, so $N\lambda=0$ in $\FF_p$.
    \item $N=2^{n}$ and $p$ is odd, so $p\nmid N$ and $\lambda=0$.
    \item Hence $\Gamma=0$, i.e.\ $[\widehat M,A]\equiv0\bmod p$.
  \end{itemize}
\end{frame}

\begin{frame}{Backup: surjectivity of reduction (Newton step)}
  Lift $g\in\Uni_N(\OK/\mm^{k})$ to an integral $g_0$ with
  $g_0^{*}g_0=\Id+p^{l}S$, $S^{*}=S$. Put $T=-S/2$ (integral, $p$ odd) and
  $g_1=g_0(\Id+p^{l}T)$. Then
  \[
    g_1^{*}g_1=\Id+p^{l}(T+S+T)+O(p^{2l})=\Id+O(p^{2l}).
  \]
  \begin{itemize}
    \item Each step doubles the precision and does not change $g_0$ modulo
      $\mm^{k}$.
    \item The limit $h$ is integral with $h^{*}h=\Id$, so $h\in G_N$ and
      $\pi_k(h)=g$.
    \item A continuous surjective homomorphism of a compact group carries Haar
      measure to the uniform measure on the finite image.
  \end{itemize}
\end{frame}

\end{document}
